% !TEX root = ../main.tex
\phantomsection
\section*{초록}
\label{sec:abstract}

금융 시스템은 대출을 통한 자본의 효율적 배분으로 광범위한 성장 동력을 창출해 왔다(Sch\"{a}r, 2021). 이와 같은 맥락에서 DeFi의 발전 역시 온체인 대출 및 레버리지 거래 활동의 정도에 결정적으로 달려 있을 것으로 기대된다. 그러나 DeFi 플랫폼은 담보 자산 집행의 어려움으로 인해 담보 기반 대출 실행에 제약을 받으며, 유형 자산 및 고용 정보를 포함하는 전통적 신용평가 모델을 직접 적용할 수 없기 때문에 무담보 신용 제공에도 한계가 있다(Sch\"{a}r, 2021; Ghosh et al., 2024).


기존 PageRank 기반 방법, 예를 들어 Adaptive Weighted PageRank(AWP) 알고리즘은 주로 거래량과 최근성(시간 감쇠)에 의존하여 지갑의 신뢰성을 평가한다(Do et al., 2023). 이러한 접근은 경제적 흐름과 활동 수준을 효과적으로 포착하지만, 과거 거래의 전수 파싱을 요구함으로써 상당한 계산 부담을 초래할 수 있다.


본 연구는 ERC-20 approve 및 permit 승인(allowance)을 지갑 보증(endorsement)의 직접적 지표로 활용하는 온체인 평판 지표 EndorseRank를 제안한다. 각 approve/permit 행위는 다른 지갑에 지출 권한을 명시적으로 부여하며, 본질적으로 신뢰를 나타낸다. 각 지갑 쌍에 대해 가장 최근 승인 금액만을 기반으로 방향성 가중 그래프 $G_{\text{approve}}$를 구성함으로써, EndorseRank는 과거 시간 감쇠 조정의 필요성을 제거하고 계산 복잡도를 크게 줄인다.


제안하는 EndorseRank 알고리즘은 \textbf{Arbitrum One}에서 이 승인 기반 그래프에 PageRank 방법론을 적용한다. 실증 검증은 \textbf{GMX V2 무기한 스왑 포지션} 결과를 사용한다: 롱·숏 포지션에 대한 담보화된 레버리지 노출이며, 수익성 있는 \textbf{PositionDecrease} 청산(양의 \textbf{basePnlUsd}, 포지션 청산 제외)은 차입 노출의 성공적 관리에 대한 기능적 유사치로 취급된다. 본 연구의 목표는 다음과 같다.


\begin{enumerate}
\item[1.] Arbitrum One에서 최신 approve/permit 승인을 추출·관리하는 견고한 데이터 파이프라인을 개발한다.
\item[2.] 희소하고 보증 지향적인 그래프로부터 평판 점수를 구성하고 효율적으로 계산한다.
\item[3.] Spearman's rho 및 Kendall's tau를 사용하여 EndorseRank가 GMX 거래 성공 대리변수(청산 성공 횟수 및 실현 이익 대리변수)와 AWP 및 기타 확립된 거래 기반 방법 대비 얼마나 정렬되는지 평가한다.
\item[4.] 런타임 및 메모리 사용량에서 우수한 효율성을 입증하기 위한 광범위한 계산 벤치마크를 수행한다.
\end{enumerate}
이러한 목표를 달성함으로써 EndorseRank는 실제 DeFi 플랫폼에 적합한 계산 효율적이고 의미론적으로 명확하며 도메인에 정렬된 평판 점수 산출 메커니즘을 제공하여, 더 나은 리스크 관리와 의사결정을 가능하게 하고자 한다.


\textbf{\emph{키워드--- }}\textbf{DeFi, 온체인 평판, GMX V2, 무기한 스왑 포지션, Arbitrum One, 블록체인 데이터 분석, 지갑 보증 그래프}
